\documentclass[12pt]{article}
\usepackage[utf8]{inputenc}
\usepackage[T1]{fontenc}
\usepackage[margin=1in]{geometry}
\usepackage{setspace}
\usepackage{xcolor}
\usepackage[colorlinks=true, urlcolor=blue!60!black, linkcolor=blue!60!black]{hyperref}
\setlength{\parindent}{0em}
\setlength{\parskip}{0.45em}
\renewcommand{\baselinestretch}{1.0}
\begin{document}
\begin{center}
{\Large \textit{Cohen v.\ California}, 403 U.S.\ 15 (1971)\par}
\vspace{0.3em}
{\normalsize Case Brief --- One-Page Edition\par}
\end{center}

\textbf{Memory Jogger.} A ``Fuck the Draft'' jacket in a courthouse corridor was protected speech: not obscene, not fighting words, and viewers could look away (\href{https://supreme.justia.com/cases/federal/us/403/15/}{Justia}).

\textbf{Facts.} In 1968, nineteen-year-old Paul Robert Cohen wore the jacket in the Los Angeles County Courthouse, where women and children were present, to protest the Vietnam War and the draft. He made no noise and provoked no violence.

\textbf{Procedural History.} Convicted of disturbing the peace by ``offensive conduct'' and sentenced to 30 days; the California Court of Appeal affirmed.

\textbf{Votes.} Reversed 5--4: Harlan for the Court; Blackmun, with Burger and Black (and White in part), dissented.

\textbf{Holding.} Absent a particularized, compelling reason, a State may not criminalize the public display of a single four-letter expletive conveying a political message.

\textbf{Reasoning.} The conviction punished words, not conduct, under a statewide statute that was no decorum rule. The slogan was not erotic, not aimed at anyone present, and avoidable ``simply by averting their eyes.'' Officials cannot draw principled lines among offensive words, emotive force deserves protection, and word-censorship risks suppressing ideas. Blackmun called it ``mainly conduct, and little speech.''

\textbf{Examples.} \emph{Same-Side:} a profane protest sign outside city hall. \emph{Opposite-Side:} a face-to-face epithet daring a fight (\textit{Chaplinsky}). \emph{Fence-Sitter:} a narrow ban on profane apparel in a courtroom.

\textbf{Critique.} ``Avert your eyes'' may undervalue public decorum, and later cases (\textit{Pacifica}, \textit{Fraser}) drew lines around profanity by setting.

\textbf{Key Quotation.} ``[I]t is nevertheless often true that one man's vulgarity is another's lyric. Indeed, we think it is largely because governmental officials cannot make principled distinctions in this area that the Constitution leaves matters of taste and style so largely to the individual.'' --- Justice Harlan, 403 U.S.\ at 25.

\end{document}
